\section{Introduction}
\label{sec:introduction}

Monitoring of the maritime environment and its associated infrastructure 
represents one of the most demanding and rapidly growing challenges in 
modern robotics. The ocean covers more than 70\% 
of the Earth's surface and hosts a vast array of critical activities, 
each requiring continuous observation and rapid response capability. 
Search and Rescue (SAR) operations must locate survivors within a narrow 
survival window that shrinks dramatically with water temperature and sea 
state \cite{iamsar2022}. Subsea pipeline networks spanning thousands of 
kilometres must be regularly inspected for leaks and structural fatigue 
to prevent environmental catastrophe. Offshore 
oil and gas platforms require around-the-clock perimeter monitoring for 
both safety and security purposes \cite{eu_blue_economy_infra2026}. Wind farm 
arrays, increasingly deployed far from shore, demand periodic structural 
inspection of turbine foundations in conditions that are often hostile to 
human operators \cite{winspector2019}.

Crewed surface vessels are expensive to procure, costly to operate, and require skilled personnel who must be exposed to dangers of the sea, especially in harsh weather conditions. A single offshore 
patrol vessel can cost several million euros per year to operate, yet 
its coverage at any given moment is limited to the vicinity of a single 
hull.

For example, single ship type of Search and Rescue Vessel 1605 FRP produced by Damen Shipyards, which is a common type of vessel used for SAR operations, has a price of around 2 million euros \cite{damen_sar1605}, depending on the configuration. Operating costs for such a vessel can easily exceed 500,000 euros per year when accounting for crew salaries, fuel, maintenance, insurance and crew training.

Manned helicopter, while faster, are 
constrained by crew endurance, weather, and the finite number of 
aircraft available at any regional rescue coordination centre. In time-
critical scenarios such as a person in the water, these constraints can 
prove fatal.

Autonomous Unmanned Surface Vehicles offer a fundamentally 
different operational model. A swarm of small, inexpensive USVs can be 
pre-positioned across a region of interest, like (big harbours, 
anchorage areas or marine work areas), and 
operated continuously without risk to human life. Because 
each vehicle in the swarm covers an independent portion of the search or 
inspection area simultaneously, total coverage time scales inversely 
with swarm size a property no single crewed platform can replicate. 
The cost per vehicle is orders of magnitude lower than that of a crewed 
vessel, enabling disposable or semi-disposable deployment strategies 
that would be economically unthinkable with conventional assets. For 
applications such as SAR, pipeline survey, platform perimeter patrol, 
and wind farm inspection, USV swarms therefore represent not merely an 
incremental improvement but a qualitative shift in what is operationally 
possible.


\subsection{Military, Critical Infrastructure and Sea Biosphere Monitoring}

The strategic importance of maritime infrastructure monitoring was 
brought into sharp focus in September 2022, when the Nord Stream 1 and 
Nord Stream 2 natural gas pipelines were destroyed in the Baltic Sea 
\cite{nordstream_website}. The incident severed 
approximately 1,200 kilometres of pipeline carrying up to 55 billion 
cubic metres of gas annually and caused one of the largest single 
methane release events ever recorded \cite{mohrmann2025nordstream}. 
Critically, the damage went undetected for hours despite the pipelines 
passing through one of the most heavily monitored bodies of water in 
the world. The event exposed a fundamental vulnerability: existing 
maritime surveillance infrastructure is ill-suited to detecting 
deliberate or accidental damage to subsea assets in a timely manner. 
Nord Stream was not an isolated case. The Baltic Sea alone hosts 
dozens of active submarine power cables and fibre-optic communication 
links connecting the national grids and internet backbones of Finland, 
Estonia, Latvia, Lithuania, Sweden, Denmark, Germany, and Poland 
\cite{balticcable}. In October 2023, the Balticconnector gas pipeline 
and a telecommunications cable between Finland and Estonia were severed, 
again under circumstances that were not immediately detected by any 
automated monitoring system \cite{cinea_balticconnector}. These incidents demonstrate that subsea critical infrastructure represents a need for continous monitoring. The economic and geopolitical consequences of such disruptions are profound, with ripple effects.

\vspace{1cm}

Beyond physical infrastructure, the maritime environment itself demands 
continuous observation. Ocean ecosystems are under sustained pressure 
from climate change, pollution, and over-exploitation.
Harmful algal blooms can develop within days 
and devastate aquaculture industries covering hundreds of square 
kilometres before they are detected by satellite or sporadic crewed 
survey. 
Illegal, Unreported, and Unregulated (IUU) 
fishing represents a further dimension of the problem. The Food and 
Agriculture Organisation of the United Nations estimates that IUU 
fishing accounts for up to 26 million tonnes of catch per year, 
representing as much as 15\% of total global marine capture 
\cite{ec_iuu_fishing2024}. Tracking fishing vessels across exclusive economic 
zones spanning hundreds of thousands of square kilometres exceeds the 
practical capacity of coast guard fleets in most nations, and vessels 
engaged in illegal activity frequently disable or spoof their AIS 
transponders to avoid detection \cite{marinetraffic_ais_spoofing}. Persistent, 
wide-area surveillance by autonomous ships capable of 
passively logging AIS traffic, detecting vessels that have gone dark, 
and collecting in-situ environmental samples offers a great 
response to both challenges simultaneously.

\vspace{1cm}

\subsection{Summary}
UAVs of course cannot replace crewed vessels and helicopters in all scenarios. But it can complement them in many, and in some cases it can enable entirely new operational concepts that were previously infeasible. For example, a swarm of USVs could be deployed to continuously monitor a busy harbour for security threats, providing real-time data to a central command centre and allowing crewed patrol boats to focus on response rather than routine surveillance. In offshore wind farms, USVs could perform regular inspections of turbine foundations, reducing the need for human divers and increasing the frequency of inspections. In SAR operations, USVs could be rapidly deployed to search large areas of open water, providing critical information to rescue teams and potentially saving lives. 




\section{Research Questions}
\label{sec:research_questions}

This thesis is motivated by the absence of a simulation platform that 
simultaneously addresses marine dynamics, swarm-scale performance, 
realistic communication constraints, and a direct path to physical 
hardware deployment. The following research questions guide the design, 
implementation, and evaluation of the system presented in this work:

\begin{description}

    \item[\textbf{RQ1} --- \textit{Simulator Fidelity and Scalability.}]
    Can a Unity3D-based simulation environment realistically model the 
    marine dynamics, virtual sensor characteristics, and LoRa radio 
    communication constraints relevant to autonomous USV swarms, while 
    maintaining real-time performance on commodity hardware?

    \item[\textbf{RQ2} --- \textit{Sensor Fusion Resilience.}]
    Can an EKF-based localisation system maintain acceptable positioning 
    accuracy for swarm coordination when individual drone GPS signals are 
    degraded or actively spoofed?

    \item[\textbf{RQ3} --- \textit{Swarm Search Performance.}]
    Can a decentralised grid-based search algorithm, coordinated 
    exclusively over LoRa, achieve complete area coverage of a realistic 
    maritime search zone within an operationally meaningful timeframe?

\end{description}

Each research question is revisited in the Conclusion 
(Section~\ref{sec:conclusion}), where the results of the experiments 
and benchmarks are used to provide a direct answer.